\section*{Chapter \ref{ch:m1:stress-case-studies} --- Stress Case Studies}

\begin{solution}{exo:m1:stress-case-studies:1}
$4.1\text{bn}/75\,000 = \$54\,667$ a contract, an index level of $54\,667/50 =
1\,093$. $35\,000/75\,000 = 47\%$ of the programme in 13 minutes, about 2\,700
contracts a minute.
\end{solution}

\begin{solution}{exo:m1:stress-case-studies:2}
At 11:00 the bands are $40 \times 0.95 = 38.00$ and $40 \times 1.05 = 42.00$; at
09:40 they are doubled, 36.00 and 44.00. An offer at 38.00 at 11:00:00 is a
limit state; if it persists the pause is declared at 11:00:15 and lasts until
11:05:15 at the earliest.
\end{solution}

\begin{solution}{exo:m1:stress-case-studies:3}
$1\,058/327 = 3.2$ halts per halted product against $220/144 = 1.5$. The funds
reopened into the same conditions that had halted them: the halt did not give
their \omterm{def:m1:what-a-trading-firm-does:market-maker}{market makers} a basket to price, so the reopening auction cleared at
the band again.
\end{solution}

\begin{solution}{exo:m1:stress-case-studies:4}
$75\,000/(0.09 \times 20\,000) = 42$ minutes; at three times the volume, 14
minutes. Volume is not liquidity: contracts passed between intermediaries
count as volume and absorb nothing. An instruction keyed to volume speeds up
exactly when its own impact provokes churn, which is when it should slow
down. A neutral instruction needs a \omterm{def:m1:asian-and-emerging-equity-markets:limit}{price limit}, a time limit, or both.
\end{solution}

\begin{solution}{exo:m1:stress-case-studies:5}
A bid of a cent and an offer of \$100\,000 are a two-sided quote, so the
obligation is met while no trade is intended. When real quotes were pulled,
the \omterm{def:m1:stress-case-studies:flash}{stub quotes} became the best prices in the book and market orders, among
them stop-loss orders that had become market orders, executed against them;
these were a large part of the trades later cancelled. In November 2010 the
regulator approved rules requiring \omterm{def:m1:what-a-trading-firm-does:market-maker}{market makers}' quotes to stay within a set
percentage of the \omterm{def:m1:us-equity-market-structure:nbbo}{national best bid and offer}, which prohibits \omterm{def:m1:stress-case-studies:flash}{stub quotes} in
effect.
\end{solution}

\begin{solution}{exo:m1:stress-case-studies:6}
The band is the cost of trading the basket plus the uncertainty on its value.
With many constituents not yet open, or open with wide and thin quotes, the
hedge could not be priced or executed: the cost-of-hedging term was tens of
per cent wide for those minutes. A pause in the fund does nothing to open the
constituents; only time did.
\end{solution}

\begin{solution}{exo:m1:stress-case-studies:7}
The default parameters give a trough of 9.8\%; both feedbacks set to zero,
2.2\%; a pause at 0.05, 7.0\%.
Churn alone: 2.2\%, exactly as with no feedback, because with a fixed book the
cost of 75\,000 contracts does not depend on how fast they are sold.
Withdrawal alone: 4.0\%. The full model at 3\% participation: 2.6\%.
Withdrawal responds to the speed of the fall; churn sets the speed of the
selling. Withdrawal without churn meets a programme slow enough for depth to
come back; churn without withdrawal accelerates a programme in a book that
does not care. Together they form the loop.
\end{solution}

\begin{solution}{exo:m1:stress-case-studies:8}
(i) \emph{Intraday inventory.} Flat at the close says nothing about 14:45; the
loss of a crash is taken and, with luck, recovered within the day, and the
risk is being stopped out, or having the winning leg cancelled, in between.
(ii) \emph{Broken trades and halts.} The backtest assumes every fill stands
and every hedge can be executed; on these days purchases were cancelled while
the hedges sold against them stood, and hedges were halted. (iii) \emph{No such
day in the sample.} Five calm years contain no withdrawal of the other \omterm{def:m1:what-a-trading-firm-does:market-maker}{market
makers}, no stale or missing data, no \omterm{def:m1:financing:margincall}{margin call}. Widening ``when volatility
rises'' reacts to a measured volatility that lags by minutes an event that
lasts minutes.
\end{solution}

\begin{solution}{pb:m1:stress-case-studies:1}
\textbf{1.} $0.09 \times 20\,000 = 1\,800$ contracts.
\textbf{2.} $3\% \times 1\,800/100\,000 = 0.054\%$.
\textbf{3.} Depth $100\,000 \times (0.1 + 0.9\,e^{-1.5}) = 30\,100$ contracts; volume
$20\,000 \times 2 = 40\,000$ a minute.
\textbf{4.} The floor is 10\,000 contracts. $0.9\,e^{-150 f} = 0.01$ gives $f =
\ln 90/150 = 3.0\%$.
\textbf{5.} $0.09 \times 40\,000 = 3\,600$ contracts, costing $3\% \times 3\,600/30\,100 =
0.36\%$ in the minute, 6.7 times the first minute, which deepens the recent
fall.
\textbf{6.} $9.8/2.2 = 4.5$ times.
\textbf{7.} $75\,000 \times 54\,667 \times (9.8\% - 2.2\%)/2 = \$156$ million.
\textbf{8.} $75\,000 \times 54\,667 \times (9.8\% - 7.0\%)/2 = \$57$ million, by the same
convention.
\textbf{9.} After the pause buyers measure the fall from the reopening price, so
depth returns to its calm level and the volume the programme follows falls
back; the remaining contracts are then sold into a full book. The report
describes the same thing: in the five seconds sell-side pressure eased,
buy-side interest returned, and the price stabilised when trading resumed.
\textbf{10.} 161\,200 contracts for 2\%; 116\,200 for 5\%.
\textbf{11.} 1.61 and 1.16 times.
\textbf{12.} With no feedback the fall is about $3\% \times 75\,000/\text{depth}$, which
is 2.25\% at 100\,000: the book was nearly deep enough. The feedbacks switch
on only when the fall outruns the five-minute average; a little more depth
keeps the early fall slow, depth then never withdraws, and the loop never
starts. Below that threshold the system is in another regime.
\textbf{13.} Sales of 600 a minute at first; the programme takes about two hours;
the price never gets far below its recent average, so neither feedback
engages: the simulation gives 2.6\%. The cost to the seller is time, that is,
two hours of exposure to the market risk the hedge was meant to remove.
\textbf{14.} The seller, or their broker: a \omterm{def:m1:asian-and-emerging-equity-markets:limit}{price limit} or a slower rate costs
almost nothing to set. Every other remedy (more capital at \omterm{def:m1:what-a-trading-firm-does:market-maker}{market makers},
exchange pauses) is more expensive and less certain.
\textbf{15.} Not in net terms. They provided immediacy to each seller for a few
seconds and passed the position on; the market's capacity to hold the 75\,000
contracts did not increase. Volume measured their activity, not absorption.
\textbf{16.} For: the prices carried no information, and the sellers were largely
individuals whose stop orders behaved in a way they had not understood.
Against: those who bought on the way down and hedged elsewhere were left with
one leg, which teaches liquidity providers to withdraw next time; and a
threshold chosen afterwards is a lottery. The second argument is why a
threshold fixed in advance is better than one chosen afterwards.
\textbf{17.} Fixed: trades at absurd prices, which the bands forbid. Not fixed:
the absence of an anchor. Funds were halted repeatedly and reopened at the
band, because a halt does not make the basket priceable.
\textbf{18.} A \omterm{def:m1:asian-and-emerging-equity-markets:limit}{price limit}, a maximum rate in contracts and not only in percentage
of volume, and a rule that pauses it and calls a human when the price has
moved more than a set amount since it started.
\textbf{19.} \textbf{The depth that was missing:} in this model, 161\,200 contracts, 61\%
more than the 100\,000 there were, would have held the fall to 2\%.
\textbf{20.} Liquidity is the quantity that someone is still willing to buy after
the price has already fallen fast, and it is smaller, by an order of
magnitude, than the quantity displayed beforehand.
\end{solution}

\begin{solution}{iq:m1:stress-case-studies:1}
A large hedger sold 75\,000 E-mini futures through an algorithm that followed
volume with no price or time limit, in a market already stressed.
Intermediaries absorbed the first sales, reached their limits and traded
among themselves, which raised volume and hence the selling rate. Depth in
the future fell to about 1\% of normal; a five-second exchange pause turned
it. In shares and funds \omterm{def:m1:what-a-trading-firm-does:market-maker}{market makers} withdrew, leaving \omterm{def:m1:stress-case-studies:flash}{stub quotes}, and
20\,000 trades at prices 60\% or more away were cancelled.
\iqlookfor{the mechanism (flow, withdrawal, feedback through volume), not just
``an algorithm did it''; the split between futures and shares.}
\end{solution}

\begin{solution}{iq:m1:stress-case-studies:2}
Bands of 5\%, 10\% or 20\% around the mean trade price of the last five
minutes, doubled near the open and the close; no trade outside them. If the
best offer sits at the lower band, or the bid at the upper, for 15 seconds,
the primary exchange pauses the security for at least five minutes on all
venues and reopens it with an auction. The reference moves only when the
mean has moved 1\%, and at most every 30 seconds.
\iqlookfor{bands on quotes versus pauses; the moving reference; the reopening
auction as the weak point.}
\end{solution}

\begin{solution}{iq:m1:stress-case-studies:3}
Their price is held to their value by arbitrageurs who must price and trade
the basket. At the open many constituents had not opened or had no reliable
quotes, futures had been limit down, and the reference values published for
the funds were stale. With no hedge, \omterm{def:m1:what-a-trading-firm-does:market-maker}{market makers} quoted wide or not at
all; market and stop orders met an empty book; halts followed, and
reopenings into the same vacuum halted again.
\iqlookfor{the arbitrage band widening, not ``panic''; why the halt did not
help.}
\end{solution}

\begin{solution}{iq:m1:stress-case-studies:4}
Not from the backtest's own distribution. Replay the strategy through
recorded or reconstructed stress days, with the things a backtest assumes
away switched off: fills that are later broken, hedges that are halted, data
that is stale or missing, latency and queue position that degrade, margin
that rises. Look at the worst intraday inventory and the capital needed at
that instant, not the daily result. Decide the stop rule in advance and test
the rule.
\iqlookfor{scenario replay with broken assumptions; intraday and not
end-of-day risk; humility about the sample.}
\end{solution}

\begin{solution}{iq:m1:stress-case-studies:5}
By rule decided in advance, not by improvisation: reduce size and widen as
inventory and short-horizon volatility rise; stop adding to inventory at its
limit; check data integrity (stale feeds, crossed markets, a hedge venue that
is halted) and stop quoting what cannot be hedged or priced. A kill switch
belongs to the trader and to risk independently. Switching off entirely is
legitimate but has a cost, including obligations if the firm is a registered
\omterm{def:m1:what-a-trading-firm-does:market-maker}{market maker}, so it should be a designed state and not a panic.
\iqlookfor{pre-defined states; the hedge and data checks; who holds the
switch; awareness of \omterm{def:m1:options-market-structure:mm}{quoting obligations}.}
\end{solution}

\begin{solution}{iq:m1:stress-case-studies:6}
A flow that does not respond to price (volume-following, stops, daily
rebalances, \omterm{def:m1:financing:margincall}{margin calls}); intermediaries with a hard limit (inventory,
capital, a hedge that must exist, doubt that trades will stand); and a loop
from the price move back to the flow or the limit. To anticipate: look for
large mechanical flows relative to the market they must trade in (leveraged
products' rebalance against futures volume, \omterm{def:m1:zero-day-options-and-retail-flow:dealer}{dealer gamma}, crowded shorts
against float), for concentration of liquidity provision in a few firms with
similar limits, and for places where the circuit breaker itself removes the
anchor.
\iqlookfor{a general structure, then concrete measurable candidates.}
\end{solution}
